\documentclass[11pt]{article}
\usepackage[margin=1.1in]{geometry}
\usepackage{amsmath,amssymb,amsthm}
\usepackage{enumitem}
\usepackage{hyperref}

\newtheorem{theorem}{Theorem}
\newtheorem{lemma}{Lemma}
\newtheorem{corollary}{Corollary}
\newtheorem{remark}{Remark}

\newcommand{\N}{\mathbb N}
\newcommand{\ellone}{\ell_1}
\newcommand{\norm}[1]{\left\lVert #1 \right\rVert}
\newcommand{\dual}[1]{#1^*}
\newcommand{\weakp}{\mu_{p,n}}

\title{Schur, WSC, and \(c_0\) Classes for Almost Multi-Boundedness}
\author{fa\_banach\_001, lane 17}
\date{2026-06-20}

\begin{document}
\maketitle

\section*{Source and Target}

The source paper is Hung Le Pham, \emph{A combinatorial characterisation of
amenable locally compact groups}, arXiv:2310.19178.  Problem 4.8 asks:
\begin{quote}
Determine those classes of Banach spaces in which relative weak compactness
implies and/or is implied by almost \((p,q)\)-multi-boundedness.
\end{quote}
The paper proves the equivalence for \(\mathfrak L^1\)-spaces and gives
examples showing that neither implication holds in arbitrary Banach spaces.
The result below gives additional positive and negative classes for every
\(1\le p,q<\infty\).

\section*{Definitions}

For a Banach space \(E\), \(x=(x_1,\ldots,x_n)\in E^n\), and
\(1\le p,q<\infty\), Pham defines
\[
  \norm{x}^{(p,q)}_n =
  \sup\left\{\left(\sum_{j=1}^n |\lambda_j(x_j)|^q\right)^{1/q}:
  \lambda_1,\ldots,\lambda_n\in E^*,\
  \sup_{\norm{u}\le 1}\left(\sum_{j=1}^n|\lambda_j(u)|^p\right)^{1/p}\le 1
  \right\}.
\]
For \(B\subset E\), write \(c_{p,q}(B)\) for the common quantity in Pham's
Proposition 3.4:
\[
  c_{p,q}(B)=
  \lim_{n\to\infty}
  \sup_{x_1,\ldots,x_n\in B}
  \frac{\norm{(x_1,\ldots,x_n)}^{(p,q)}_n}{n^{1/q}}.
\]
The set \(B\) is almost \((p,q)\)-multi-bounded when \(c_{p,q}(B)=0\).

\section*{Proof Intuition}

The tuple norms are strong enough to detect separated \(\ell_1\)-structure,
but weak enough that finite sets disappear after normalization by \(n^{1/q}\).
A norm-compact set is approximated by a finite set, so it is automatically
almost multi-bounded.  For the implication from almost multi-boundedness to
weak compactness, Rosenthal's dichotomy is decisive: in weakly sequentially
complete spaces a bounded non-weakly compact set must contain an \(\ell_1\)
basic sequence, and such a sequence is seen by the tuple norms.  On the other
hand, the summing basis in \(c_0\), already isolated in Pham's examples, gives
an almost multi-bounded set that is not weakly compact; this obstruction passes
to every space containing \(c_0\).

\section*{Lemmas}

\begin{lemma}[Finite sets vanish asymptotically]
\label{lem:finite}
Let \(A\) be a finite subset of a Banach space \(E\).  Then
\(c_{p,q}(A)=0\) for every \(1\le p,q<\infty\).
\end{lemma}

\begin{proof}
Write \(A=\{a_1,\ldots,a_m\}\).  Let \(x_1,\ldots,x_n\in A\), and for each
\(r\) put \(I_r=\{j:x_j=a_r\}\).  If
\(\lambda_1,\ldots,\lambda_n\in E^*\) satisfy
\[
  \sup_{\norm{u}\le 1}\left(\sum_{j=1}^n|\lambda_j(u)|^p\right)^{1/p}\le 1,
\]
then
\[
  \left(\sum_{j\in I_r}|\lambda_j(a_r)|^p\right)^{1/p}\le \norm{a_r}.
\]
If \(q\ge p\), the corresponding \(q\)-norm is at most \(\norm{a_r}\).  If
\(q<p\), it is at most
\(n^{1/q-1/p}\norm{a_r}\).  Hence
\[
  \norm{(x_1,\ldots,x_n)}^{(p,q)}_n
  \le
  \left(\sum_{r=1}^m \norm{a_r}\right)n^{\max\{0,1/q-1/p\}}.
\]
After division by \(n^{1/q}\), this tends to \(0\).
\end{proof}

\begin{lemma}[Norm-compact sets are almost multi-bounded]
\label{lem:compact}
Let \(B\) be norm relatively compact in a Banach space \(E\).  Then
\(c_{p,q}(B)=0\) for every \(1\le p,q<\infty\).
\end{lemma}

\begin{proof}
Fix \(\eta>0\), and choose a finite \(\eta\)-net \(A\) for \(B\).  For
\(x_1,\ldots,x_n\in B\), choose \(a_j\in A\) with \(\norm{x_j-a_j}\le \eta\).
Every admissible tuple \((\lambda_j)\) in the definition of
\(\norm{\cdot}^{(p,q)}_n\) satisfies \(\norm{\lambda_j}\le 1\) for each \(j\).
Thus
\[
  \norm{(x_1-a_1,\ldots,x_n-a_n)}^{(p,q)}_n\le \eta n^{1/q}.
\]
By Lemma \ref{lem:finite},
\[
  \limsup_{n\to\infty}\sup_{x_1,\ldots,x_n\in B}
  \frac{\norm{(x_1,\ldots,x_n)}^{(p,q)}_n}{n^{1/q}}\le \eta.
\]
Letting \(\eta\downarrow 0\) proves the claim.
\end{proof}

\begin{lemma}[\(\ell_1\)-basic obstruction]
\label{lem:l1}
If \(B\subset E\) contains a sequence equivalent to the unit vector basis of
\(\ell_1\), then \(B\) is not almost \((p,q)\)-multi-bounded for any
\(1\le p,q<\infty\).
\end{lemma}

\begin{proof}
Pass to a sequence \((x_j)\subset B\) and choose \(\kappa>0\) such that
\[
  \norm{\sum_{j=1}^n \alpha_j x_j}\ge
  \kappa\sum_{j=1}^n |\alpha_j|
\]
for all finite scalar families \((\alpha_j)\).  The coordinate map
\(\sum \alpha_jx_j\mapsto(\alpha_j)\) from \([x_j]\) into \(\ell_1\) has norm
at most \(\kappa^{-1}\).  Let \(\lambda_j\in E^*\) be Hahn--Banach extensions
of its coordinates.  For every \(u\in E\),
\[
  \left(\sum_{j=1}^n |\lambda_j(u)|^p\right)^{1/p}
  \le \sum_{j=1}^n |\lambda_j(u)|
  \le \kappa^{-1}\norm{u}.
\]
Therefore \((\kappa\lambda_1,\ldots,\kappa\lambda_n)\) is admissible in the
definition of the \((p,q)\)-multi-norm, and
\[
  \norm{(x_1,\ldots,x_n)}^{(p,q)}_n
  \ge
  \left(\sum_{j=1}^n|\kappa\lambda_j(x_j)|^q\right)^{1/q}
  =\kappa n^{1/q}.
\]
Hence \(c_{p,q}(B)\ge \kappa>0\).
\end{proof}

\begin{lemma}[Functoriality]
\label{lem:operator}
Let \(T:E\to F\) be a bounded linear operator.  If \(B\subset E\) is almost
\((p,q)\)-multi-bounded, then \(T(B)\subset F\) is almost
\((p,q)\)-multi-bounded.
\end{lemma}

\begin{proof}
For every \(x_1,\ldots,x_n\in E\) and every admissible tuple
\((\lambda_1,\ldots,\lambda_n)\in (F^*)^n\), the tuple
\((T^*\lambda_1,\ldots,T^*\lambda_n)\) has weak \(p\)-summing norm at most
\(\norm{T}\).  Therefore
\[
  \norm{(Tx_1,\ldots,Tx_n)}^{(p,q)}_n
  \le \norm{T}\,\norm{(x_1,\ldots,x_n)}^{(p,q)}_n.
\]
After taking suprema and normalizing by \(n^{1/q}\), the assertion follows.
\end{proof}

\section*{Theorem}

\begin{theorem}
\label{thm:main}
Let \(1\le p,q<\infty\).
\begin{enumerate}[label=\rm(\alph*)]
\item In every Banach space, norm relatively compact subsets are almost
\((p,q)\)-multi-bounded.
\item If \(E\) has the Schur property, then a subset \(B\subset E\) is
relatively weakly compact if and only if it is almost
\((p,q)\)-multi-bounded.
\item If \(E\) is weakly sequentially complete, then every almost
\((p,q)\)-multi-bounded subset of \(E\) is relatively weakly compact.
\item If \(E\) contains an isomorphic copy of \(c_0\), then the implication
``almost \((p,q)\)-multi-bounded implies relatively weakly compact'' fails in
\(E\).
\end{enumerate}
\end{theorem}

\begin{proof}
Part (a) is Lemma \ref{lem:compact}.

Assume now that \(E\) has the Schur property.  If \(B\) is relatively weakly
compact, then \(B\) is norm relatively compact by Eberlein--Smulian and the
Schur property; part (a) gives almost \((p,q)\)-multi-boundedness.

Conversely, suppose \(B\) is almost \((p,q)\)-multi-bounded.  Pham observes
after Definition 3.5 that unbounded sets have \(c_{p,q}(B)=\infty\), so \(B\)
is bounded.  Since Schur spaces are weakly sequentially complete, the argument
in the next paragraph shows that \(B\) is relatively weakly compact.  This
proves the equivalence in Schur spaces.

Assume more generally that \(E\) is weakly sequentially complete, and let
\(B\) be almost \((p,q)\)-multi-bounded.  Again \(B\) is bounded.  If \(B\)
were not relatively weakly compact, Eberlein--Smulian would give a sequence in
\(B\) with no weakly convergent subsequence.  Rosenthal's \(\ell_1\) theorem
gives either an \(\ell_1\)-basic subsequence or a weakly Cauchy subsequence.
The first possibility contradicts Lemma \ref{lem:l1}.  The second contradicts
weak sequential completeness.  Hence \(B\) is relatively weakly compact.

Finally suppose that \(E\) contains a closed subspace \(Y\) isomorphic to
\(c_0\).  In the examples following Theorem 4.4, Pham records that the summing
basis
\[
  s_n=e_1+\cdots+e_n\qquad(n\in\N)
\]
forms a subset \(S=\{s_n:n\in\N\}\) of \(c_0\) that is almost
\((p,q)\)-multi-bounded for every \(p,q\), but is not relatively weakly
compact.  If \(J:c_0\to Y\subset E\) is an isomorphism, Lemma
\ref{lem:operator} shows that \(J(S)\) is almost \((p,q)\)-multi-bounded in
\(E\).  It is not relatively weakly compact: weak compactness of \(J(S)\) in
\(E\) would imply weak compactness in the closed subspace \(Y\), and then weak
compactness of \(S\) in \(c_0\) after applying \(J^{-1}\).
\end{proof}

\begin{corollary}
Every finite-dimensional Banach space has the full equivalence:
\[
  B\text{ is relatively weakly compact}
  \quad\Longleftrightarrow\quad
  B\text{ is almost }(p,q)\text{-multi-bounded}.
\]
\end{corollary}

\begin{proof}
Finite-dimensional spaces are reflexive, hence weakly sequentially complete,
and have the Schur property.
\end{proof}

\begin{corollary}
Every reflexive Banach space satisfies the implication from almost
\((p,q)\)-multi-boundedness to relative weak compactness.
\end{corollary}

\begin{proof}
Reflexive spaces are weakly sequentially complete.
\end{proof}

\section*{Verification Notes and Limitations}

This is a partial answer to Problem 4.8, not a classification.  It adds exact
families to the landscape: Schur spaces for the two-sided equivalence, weakly
sequentially complete spaces for the almost-implies-weak-compactness direction,
and spaces containing \(c_0\) as a broad negative class for that same
direction.  It is consistent with Pham's examples: infinite-dimensional
reflexive unit balls can fail to be almost multi-bounded for some parameters,
so the reverse implication does not follow from reflexivity alone.  The
remaining frontier for the almost-implies-weak-compactness direction includes
the classical non-weakly-sequentially-complete spaces that do not contain
\(c_0\), where Rosenthal's strongly summing alternative can occur.

Bounded novelty search on 2026-06-20 checked exact phrases around almost
\((p,q)\)-multi-boundedness, Schur spaces, weak sequential completeness,
\(c_0\), weak compactness, \(c_{p,q}(B)\), and arXiv:2310.19178.  No exact
existing Schur/WSC/\(c_0\) statement was found.  The result should be reviewed
as a partial landscape theorem rather than as a full solution of the source
problem.

\section*{References}

\begin{thebibliography}{9}
\bibitem{Pham}
H. L. Pham,
\emph{A combinatorial characterisation of amenable locally compact groups},
arXiv:2310.19178.

\bibitem{Rosenthal}
H. P. Rosenthal,
\emph{A characterization of Banach spaces containing \(\ell^1\)},
Proceedings of the National Academy of Sciences of the United States of America
71 (1974), 2411--2413.

\bibitem{Diestel}
J. Diestel,
\emph{Sequences and Series in Banach Spaces},
Graduate Texts in Mathematics 92, Springer, 1984.
\end{thebibliography}

\end{document}
